\documentclass[11pt,a4paper]{article}
\usepackage[margin=2.5cm]{geometry}
\usepackage[T1]{fontenc}
\usepackage{lmodern}
\usepackage{microtype}
\usepackage{booktabs}
\usepackage{amsmath,amssymb}
\usepackage{listings}
\usepackage{xcolor}
\usepackage{isabelle,isabellesym}
\isabellestyle{it}
\renewcommand{\isachardoublequote}{\textquotedbl}
\renewcommand{\isachardoublequoteopen}{\textquotedbl}
\renewcommand{\isachardoublequoteclose}{\textquotedbl}
\usepackage[hidelinks]{hyperref}
% Snippets of the Isabelle theories, extracted by make-snippets.sh
\newcommand{\usesnip}[1]{\csname #1\endcsname}
\newcommand{\snip}[4]{\expandafter\newcommand\csname #1\endcsname{#4}}
\snip{ast_ty}{1}{2}{%
\isakeywordONE{datatype}\isamarkupfalse%
\ ty\ {\isacharequal}{\kern0pt}\ TInt\ {\isacharbar}{\kern0pt}\ TBool\ {\isacharbar}{\kern0pt}\ TStr\ {\isacharbar}{\kern0pt}\ TNone\ {\isacharbar}{\kern0pt}\ TStruct\ name\ {\isacharbar}{\kern0pt}\ TList\ ty%
}%endsnip
\snip{ast_conv}{1}{2}{%
\isakeywordONE{datatype}\isamarkupfalse%
\ conv\ {\isacharequal}{\kern0pt}\ Imm\ {\isacharbar}{\kern0pt}\ Mut\ {\isacharbar}{\kern0pt}\ Var\ {\isacharbar}{\kern0pt}\ Ref\ {\isacharbar}{\kern0pt}\ Out\ {\isacharbar}{\kern0pt}\ Deinit%
}%endsnip
\snip{ast_binop}{1}{2}{%
\isakeywordONE{datatype}\isamarkupfalse%
\ binop\ {\isacharequal}{\kern0pt}\ Add\ {\isacharbar}{\kern0pt}\ Sub\ {\isacharbar}{\kern0pt}\ Mul\ {\isacharbar}{\kern0pt}\ Lt\ {\isacharbar}{\kern0pt}\ Le\ {\isacharbar}{\kern0pt}\ Eq\ {\isacharbar}{\kern0pt}\ And\ {\isacharbar}{\kern0pt}\ Or%
}%endsnip
\snip{ast_placeexpr}{1}{2}{%
\isakeywordONE{datatype}\isamarkupfalse%
\ place\ {\isacharequal}{\kern0pt}\isanewline
\ \ \ \ PVar\ vname\ %
\hfill\isamarkupcmt{\isa{x}%
}\isanewline
\ \ {\isacharbar}{\kern0pt}\ PField\ place\ fname\ %
\hfill\isamarkupcmt{\isa{p{\isachardot}{\kern0pt}f}%
}\isanewline
\ \ {\isacharbar}{\kern0pt}\ PIndex\ place\ expr\ %
\hfill\isamarkupcmt{\isa{p{\isacharbrackleft}{\kern0pt}i{\isacharbrackright}{\kern0pt}}%
}\isanewline
\isakeywordTWO{and}\ expr\ {\isacharequal}{\kern0pt}\isanewline
\ \ \ \ EInt\ int\ %
\hfill\isamarkupcmt{\isa{{\isadigit{4}}{\isadigit{2}}}%
}\isanewline
\ \ {\isacharbar}{\kern0pt}\ EBool\ bool\ %
\hfill\isamarkupcmt{\isa{True}%
}\isanewline
\ \ {\isacharbar}{\kern0pt}\ EStr\ string\ %
\hfill\isamarkupcmt{\isa{{\isachardoublequote}{\kern0pt}hi{\isachardoublequote}{\kern0pt}\ \ {\isacharparenleft}{\kern0pt}only\ as\ a\ print\ argument{\isacharparenright}{\kern0pt}}%
}\isanewline
\ \ {\isacharbar}{\kern0pt}\ EPlace\ place\ %
\hfill\isamarkupcmt{\isa{x\ \ {\isacharparenleft}{\kern0pt}a\ place\ used\ as\ a\ value{\isacharcomma}{\kern0pt}\ or\ as\ a\ reference\ argument{\isacharparenright}{\kern0pt}}%
}\isanewline
\ \ {\isacharbar}{\kern0pt}\ ETake\ place\ %
\hfill\isamarkupcmt{\isa{p{\isacharcircum}{\kern0pt}\ \ {\isacharparenleft}{\kern0pt}transfer{\isacharparenright}{\kern0pt}}%
}\isanewline
\ \ {\isacharbar}{\kern0pt}\ ECopy\ place\ %
\hfill\isamarkupcmt{\isa{p{\isachardot}{\kern0pt}copy{\isacharparenleft}{\kern0pt}{\isacharparenright}{\kern0pt}}%
}\isanewline
\ \ {\isacharbar}{\kern0pt}\ ECall\ name\ {\isachardoublequoteopen}expr\ list{\isachardoublequoteclose}\ %
\hfill\isamarkupcmt{\isa{f{\isacharparenleft}{\kern0pt}a{\isacharcomma}{\kern0pt}\ b{\isacharparenright}{\kern0pt}\ \ or\ \ Res{\isacharparenleft}{\kern0pt}{\isadigit{1}}{\isacharparenright}{\kern0pt}\ \ {\isacharparenleft}{\kern0pt}call\ or\ construction{\isacharparenright}{\kern0pt}}%
}\isanewline
\ \ {\isacharbar}{\kern0pt}\ EMeth\ place\ name\ {\isachardoublequoteopen}expr\ list{\isachardoublequoteclose}\ %
\hfill\isamarkupcmt{\isa{p{\isachardot}{\kern0pt}m{\isacharparenleft}{\kern0pt}a{\isacharparenright}{\kern0pt}\ \ {\isacharparenleft}{\kern0pt}method\ call{\isacharparenright}{\kern0pt}}%
}\isanewline
\ \ {\isacharbar}{\kern0pt}\ EBin\ binop\ expr\ expr\ %
\hfill\isamarkupcmt{\isa{a\ {\isacharplus}{\kern0pt}\ b}%
}\isanewline
\ \ {\isacharbar}{\kern0pt}\ ENot\ expr\ %
\hfill\isamarkupcmt{\isa{not\ e}%
}\isanewline
\ \ {\isacharbar}{\kern0pt}\ ELen\ place\ %
\hfill\isamarkupcmt{\isa{len{\isacharparenleft}{\kern0pt}p{\isacharparenright}{\kern0pt}}%
}\isanewline
\ \ {\isacharbar}{\kern0pt}\ EList\ {\isachardoublequoteopen}expr\ list{\isachardoublequoteclose}\ %
\hfill\isamarkupcmt{\isa{{\isacharbrackleft}{\kern0pt}{\isadigit{1}}{\isacharcomma}{\kern0pt}\ {\isadigit{2}}{\isacharcomma}{\kern0pt}\ {\isadigit{3}}{\isacharbrackright}{\kern0pt}}%
}%
}%endsnip
\snip{ast_stmt}{1}{2}{%
\isakeywordONE{datatype}\isamarkupfalse%
\ stmt\ {\isacharequal}{\kern0pt}\isanewline
\ \ \ \ SPass\ %
\hfill\isamarkupcmt{\isa{pass}%
}\isanewline
\ \ {\isacharbar}{\kern0pt}\ SSeq\ stmt\ stmt\ %
\hfill\isamarkupcmt{\isa{s{\isadigit{1}}\ followed\ by\ s{\isadigit{2}}}%
}\isanewline
\ \ {\isacharbar}{\kern0pt}\ SVar\ vname\ ty\ {\isachardoublequoteopen}expr\ option{\isachardoublequoteclose}\ %
\hfill\isamarkupcmt{\isa{var\ x{\isacharcolon}{\kern0pt}\ T\ {\isacharequal}{\kern0pt}\ e\ \ or\ \ var\ x{\isacharcolon}{\kern0pt}\ T}%
}\isanewline
\ \ {\isacharbar}{\kern0pt}\ SRef\ vname\ expr\ %
\hfill\isamarkupcmt{\isa{ref\ r\ {\isacharequal}{\kern0pt}\ e}%
}\isanewline
\ \ {\isacharbar}{\kern0pt}\ SAssign\ place\ expr\ %
\hfill\isamarkupcmt{\isa{p\ {\isacharequal}{\kern0pt}\ e}%
}\isanewline
\ \ {\isacharbar}{\kern0pt}\ SExpr\ expr\ %
\hfill\isamarkupcmt{\isa{f{\isacharparenleft}{\kern0pt}x{\isacharparenright}{\kern0pt}\ \ {\isacharparenleft}{\kern0pt}result\ discarded{\isacharparenright}{\kern0pt}}%
}\isanewline
\ \ {\isacharbar}{\kern0pt}\ SPrint\ {\isachardoublequoteopen}expr\ list{\isachardoublequoteclose}\ %
\hfill\isamarkupcmt{\isa{print{\isacharparenleft}{\kern0pt}a{\isacharcomma}{\kern0pt}\ b{\isacharparenright}{\kern0pt}}%
}\isanewline
\ \ {\isacharbar}{\kern0pt}\ SIf\ expr\ stmt\ stmt\ %
\hfill\isamarkupcmt{\isa{if\ c{\isacharcolon}{\kern0pt}\ s{\isadigit{1}}\ else{\isacharcolon}{\kern0pt}\ s{\isadigit{2}}}%
}\isanewline
\ \ {\isacharbar}{\kern0pt}\ SWhile\ expr\ stmt\ %
\hfill\isamarkupcmt{\isa{while\ c{\isacharcolon}{\kern0pt}\ s}%
}\isanewline
\ \ {\isacharbar}{\kern0pt}\ SReturn\ {\isachardoublequoteopen}expr\ option{\isachardoublequoteclose}\ %
\hfill\isamarkupcmt{\isa{return\ e\ \ or\ \ return}%
}\isanewline
\ \ {\isacharbar}{\kern0pt}\ SDrop\ place\ %
\hfill\isamarkupcmt{inserted by ASAP elaboration%
}%
}%endsnip
\snip{ast_ret}{1}{2}{%
\isakeywordONE{datatype}\isamarkupfalse%
\ ret\ {\isacharequal}{\kern0pt}\ RNone\ {\isacharbar}{\kern0pt}\ RVal\ ty\ {\isacharbar}{\kern0pt}\ RRef\ {\isachardoublequoteopen}vname\ list{\isachardoublequoteclose}\ ty%
}%endsnip
\snip{ast_fundecl}{1}{2}{%
\isakeywordONE{datatype}\isamarkupfalse%
\ fundecl\ {\isacharequal}{\kern0pt}\ Fun\isanewline
\ \ {\isacharparenleft}{\kern0pt}fd{\isacharunderscore}{\kern0pt}name{\isacharcolon}{\kern0pt}\ name{\isacharparenright}{\kern0pt}\isanewline
\ \ {\isacharparenleft}{\kern0pt}fd{\isacharunderscore}{\kern0pt}params{\isacharcolon}{\kern0pt}\ {\isachardoublequoteopen}{\isacharparenleft}{\kern0pt}conv\ {\isasymtimes}\ vname\ {\isasymtimes}\ ty{\isacharparenright}{\kern0pt}\ list{\isachardoublequoteclose}{\isacharparenright}{\kern0pt}\isanewline
\ \ {\isacharparenleft}{\kern0pt}fd{\isacharunderscore}{\kern0pt}ret{\isacharcolon}{\kern0pt}\ ret{\isacharparenright}{\kern0pt}\isanewline
\ \ {\isacharparenleft}{\kern0pt}fd{\isacharunderscore}{\kern0pt}body{\isacharcolon}{\kern0pt}\ stmt{\isacharparenright}{\kern0pt}%
}%endsnip
\snip{ast_structdecl}{1}{2}{%
\isakeywordONE{datatype}\isamarkupfalse%
\ structdecl\ {\isacharequal}{\kern0pt}\ Struct\isanewline
\ \ {\isacharparenleft}{\kern0pt}sd{\isacharunderscore}{\kern0pt}name{\isacharcolon}{\kern0pt}\ name{\isacharparenright}{\kern0pt}\isanewline
\ \ {\isacharparenleft}{\kern0pt}sd{\isacharunderscore}{\kern0pt}fields{\isacharcolon}{\kern0pt}\ {\isachardoublequoteopen}{\isacharparenleft}{\kern0pt}fname\ {\isasymtimes}\ ty{\isacharparenright}{\kern0pt}\ list{\isachardoublequoteclose}{\isacharparenright}{\kern0pt}\isanewline
\ \ {\isacharparenleft}{\kern0pt}sd{\isacharunderscore}{\kern0pt}copyable{\isacharcolon}{\kern0pt}\ bool{\isacharparenright}{\kern0pt}\isanewline
\ \ {\isacharparenleft}{\kern0pt}sd{\isacharunderscore}{\kern0pt}fieldwise{\isacharcolon}{\kern0pt}\ bool{\isacharparenright}{\kern0pt}\isanewline
\ \ {\isacharparenleft}{\kern0pt}sd{\isacharunderscore}{\kern0pt}init{\isacharcolon}{\kern0pt}\ {\isachardoublequoteopen}fundecl\ option{\isachardoublequoteclose}{\isacharparenright}{\kern0pt}\isanewline
\ \ {\isacharparenleft}{\kern0pt}sd{\isacharunderscore}{\kern0pt}deinit{\isacharcolon}{\kern0pt}\ {\isachardoublequoteopen}stmt\ option{\isachardoublequoteclose}{\isacharparenright}{\kern0pt}\isanewline
\ \ {\isacharparenleft}{\kern0pt}sd{\isacharunderscore}{\kern0pt}methods{\isacharcolon}{\kern0pt}\ {\isachardoublequoteopen}fundecl\ list{\isachardoublequoteclose}{\isacharparenright}{\kern0pt}%
}%endsnip
\snip{ast_prog}{1}{2}{%
\isakeywordONE{datatype}\isamarkupfalse%
\ prog\ {\isacharequal}{\kern0pt}\ Prog\ {\isacharparenleft}{\kern0pt}p{\isacharunderscore}{\kern0pt}structs{\isacharcolon}{\kern0pt}\ {\isachardoublequoteopen}structdecl\ list{\isachardoublequoteclose}{\isacharparenright}{\kern0pt}\ {\isacharparenleft}{\kern0pt}p{\isacharunderscore}{\kern0pt}funs{\isacharcolon}{\kern0pt}\ {\isachardoublequoteopen}fundecl\ list{\isachardoublequoteclose}{\isacharparenright}{\kern0pt}%
}%endsnip
\snip{ex_res_struct}{1}{2}{%
\isakeywordONE{definition}\isamarkupfalse%
\ res{\isacharunderscore}{\kern0pt}struct\ {\isacharcolon}{\kern0pt}{\isacharcolon}{\kern0pt}\ structdecl\ \isakeywordTWO{where}\isanewline
\ \ {\isachardoublequoteopen}res{\isacharunderscore}{\kern0pt}struct\ {\isacharequal}{\kern0pt}\ Struct\ {\isacharprime}{\kern0pt}{\isacharprime}{\kern0pt}Res{\isacharprime}{\kern0pt}{\isacharprime}{\kern0pt}\ {\isacharbrackleft}{\kern0pt}{\isacharparenleft}{\kern0pt}{\isacharprime}{\kern0pt}{\isacharprime}{\kern0pt}id{\isacharprime}{\kern0pt}{\isacharprime}{\kern0pt}{\isacharcomma}{\kern0pt}\ TInt{\isacharparenright}{\kern0pt}{\isacharbrackright}{\kern0pt}\ False\ True\ None\isanewline
\ \ \ \ \ {\isacharparenleft}{\kern0pt}Some\ {\isacharparenleft}{\kern0pt}PRINT\ {\isacharbrackleft}{\kern0pt}Txt\ {\isacharprime}{\kern0pt}{\isacharprime}{\kern0pt}drop{\isacharprime}{\kern0pt}{\isacharprime}{\kern0pt}{\isacharcomma}{\kern0pt}\ {\isasymlangle}{\isasymsection}{\isacharprime}{\kern0pt}{\isacharprime}{\kern0pt}self{\isacharprime}{\kern0pt}{\isacharprime}{\kern0pt}{\isasymbullet}{\isacharprime}{\kern0pt}{\isacharprime}{\kern0pt}id{\isacharprime}{\kern0pt}{\isacharprime}{\kern0pt}{\isasymrangle}{\isacharbrackright}{\kern0pt}{\isacharparenright}{\kern0pt}{\isacharparenright}{\kern0pt}\ {\isacharbrackleft}{\kern0pt}{\isacharbrackright}{\kern0pt}{\isachardoublequoteclose}%
}%endsnip
\snip{ex_ex1_main}{1}{2}{%
\isakeywordONE{definition}\isamarkupfalse%
\ ex{\isadigit{1}}{\isacharunderscore}{\kern0pt}main\ {\isacharcolon}{\kern0pt}{\isacharcolon}{\kern0pt}\ stmt\ \isakeywordTWO{where}\isanewline
\ \ {\isachardoublequoteopen}ex{\isadigit{1}}{\isacharunderscore}{\kern0pt}main\ {\isacharequal}{\kern0pt}\isanewline
\ \ \ \ \ {\isacharparenleft}{\kern0pt}VAR\ {\isacharprime}{\kern0pt}{\isacharprime}{\kern0pt}a{\isacharprime}{\kern0pt}{\isacharprime}{\kern0pt}\ {\isasymZtypecolon}\ Res\ {\isacharcolon}{\kern0pt}{\isacharcolon}{\kern0pt}{\isacharequal}{\kern0pt}\ {\isacharprime}{\kern0pt}{\isacharprime}{\kern0pt}Res{\isacharprime}{\kern0pt}{\isacharprime}{\kern0pt}{\isasymlbrace}{\isacharbrackleft}{\kern0pt}N\ {\isadigit{1}}{\isacharbrackright}{\kern0pt}{\isasymrbrace}\ {\isacharsemicolon}{\kern0pt}{\isacharsemicolon}{\kern0pt}\isanewline
\ \ \ \ \ VAR\ {\isacharprime}{\kern0pt}{\isacharprime}{\kern0pt}b{\isacharprime}{\kern0pt}{\isacharprime}{\kern0pt}\ {\isasymZtypecolon}\ Res\ {\isacharcolon}{\kern0pt}{\isacharcolon}{\kern0pt}{\isacharequal}{\kern0pt}\ {\isacharprime}{\kern0pt}{\isacharprime}{\kern0pt}Res{\isacharprime}{\kern0pt}{\isacharprime}{\kern0pt}{\isasymlbrace}{\isacharbrackleft}{\kern0pt}N\ {\isadigit{2}}{\isacharbrackright}{\kern0pt}{\isasymrbrace}\ {\isacharsemicolon}{\kern0pt}{\isacharsemicolon}{\kern0pt}\isanewline
\ \ \ \ \ EVAL\ {\isacharprime}{\kern0pt}{\isacharprime}{\kern0pt}bump{\isacharprime}{\kern0pt}{\isacharprime}{\kern0pt}{\isasymlbrace}{\isacharbrackleft}{\kern0pt}{\isasymlangle}{\isasymsection}{\isacharprime}{\kern0pt}{\isacharprime}{\kern0pt}a{\isacharprime}{\kern0pt}{\isacharprime}{\kern0pt}{\isasymrangle}{\isacharcomma}{\kern0pt}\ N\ {\isadigit{1}}{\isadigit{0}}{\isacharbrackright}{\kern0pt}{\isasymrbrace}\ {\isacharsemicolon}{\kern0pt}{\isacharsemicolon}{\kern0pt}\isanewline
\ \ \ \ \ REF\ {\isacharprime}{\kern0pt}{\isacharprime}{\kern0pt}r{\isacharprime}{\kern0pt}{\isacharprime}{\kern0pt}\ {\isacharcolon}{\kern0pt}{\isacharcolon}{\kern0pt}{\isacharequal}{\kern0pt}\ {\isacharprime}{\kern0pt}{\isacharprime}{\kern0pt}pick{\isacharprime}{\kern0pt}{\isacharprime}{\kern0pt}{\isasymlbrace}{\isacharbrackleft}{\kern0pt}TT{\isacharcomma}{\kern0pt}\ {\isasymlangle}{\isasymsection}{\isacharprime}{\kern0pt}{\isacharprime}{\kern0pt}a{\isacharprime}{\kern0pt}{\isacharprime}{\kern0pt}{\isasymrangle}{\isacharcomma}{\kern0pt}\ {\isasymlangle}{\isasymsection}{\isacharprime}{\kern0pt}{\isacharprime}{\kern0pt}b{\isacharprime}{\kern0pt}{\isacharprime}{\kern0pt}{\isasymrangle}{\isacharbrackright}{\kern0pt}{\isasymrbrace}\ {\isacharsemicolon}{\kern0pt}{\isacharsemicolon}{\kern0pt}\isanewline
\ \ \ \ \ PRINT\ {\isacharbrackleft}{\kern0pt}{\isasymlangle}{\isasymsection}{\isacharprime}{\kern0pt}{\isacharprime}{\kern0pt}r{\isacharprime}{\kern0pt}{\isacharprime}{\kern0pt}{\isasymbullet}{\isacharprime}{\kern0pt}{\isacharprime}{\kern0pt}id{\isacharprime}{\kern0pt}{\isacharprime}{\kern0pt}{\isasymrangle}{\isacharbrackright}{\kern0pt}\ {\isacharsemicolon}{\kern0pt}{\isacharsemicolon}{\kern0pt}\isanewline
\ \ \ \ \ EVAL\ {\isacharprime}{\kern0pt}{\isacharprime}{\kern0pt}consume{\isacharprime}{\kern0pt}{\isacharprime}{\kern0pt}{\isasymlbrace}{\isacharbrackleft}{\kern0pt}{\isasymsection}{\isacharprime}{\kern0pt}{\isacharprime}{\kern0pt}a{\isacharprime}{\kern0pt}{\isacharprime}{\kern0pt}\isactrlsup {\isacharcircum}{\kern0pt}{\isacharbrackright}{\kern0pt}{\isasymrbrace}\ {\isacharsemicolon}{\kern0pt}{\isacharsemicolon}{\kern0pt}\isanewline
\ \ \ \ \ VAR\ {\isacharprime}{\kern0pt}{\isacharprime}{\kern0pt}l{\isacharprime}{\kern0pt}{\isacharprime}{\kern0pt}\ {\isasymZtypecolon}\ TList\ TInt\ {\isacharcolon}{\kern0pt}{\isacharcolon}{\kern0pt}{\isacharequal}{\kern0pt}\ LIST\ {\isacharbrackleft}{\kern0pt}N\ {\isadigit{1}}{\isacharcomma}{\kern0pt}\ N\ {\isadigit{2}}{\isacharcomma}{\kern0pt}\ N\ {\isadigit{3}}{\isacharbrackright}{\kern0pt}\ {\isacharsemicolon}{\kern0pt}{\isacharsemicolon}{\kern0pt}\isanewline
\ \ \ \ \ REF\ {\isacharprime}{\kern0pt}{\isacharprime}{\kern0pt}e{\isacharprime}{\kern0pt}{\isacharprime}{\kern0pt}\ {\isacharcolon}{\kern0pt}{\isacharcolon}{\kern0pt}{\isacharequal}{\kern0pt}\ {\isasymlangle}{\isasymsection}{\isacharprime}{\kern0pt}{\isacharprime}{\kern0pt}l{\isacharprime}{\kern0pt}{\isacharprime}{\kern0pt}{\isasymlbrakk}N\ {\isadigit{0}}{\isasymrbrakk}{\isasymrangle}\ {\isacharsemicolon}{\kern0pt}{\isacharsemicolon}{\kern0pt}\isanewline
\ \ \ \ \ PRINT\ {\isacharbrackleft}{\kern0pt}{\isasymlangle}{\isasymsection}{\isacharprime}{\kern0pt}{\isacharprime}{\kern0pt}e{\isacharprime}{\kern0pt}{\isacharprime}{\kern0pt}{\isasymrangle}{\isacharbrackright}{\kern0pt}\ {\isacharsemicolon}{\kern0pt}{\isacharsemicolon}{\kern0pt}\isanewline
\ \ \ \ \ EVAL\ {\isasymsection}{\isacharprime}{\kern0pt}{\isacharprime}{\kern0pt}l{\isacharprime}{\kern0pt}{\isacharprime}{\kern0pt}{\isasymbullet}{\isacharprime}{\kern0pt}{\isacharprime}{\kern0pt}append{\isacharprime}{\kern0pt}{\isacharprime}{\kern0pt}{\isasymlbrace}{\isacharbrackleft}{\kern0pt}N\ {\isadigit{4}}{\isacharbrackright}{\kern0pt}{\isasymrbrace}{\isacharparenright}{\kern0pt}{\isachardoublequoteclose}%
}%endsnip

\snip{invoke_eqn}{1}{2}{%
%
\begin{isamarkuptext}%
%
\begin{isabelle}%
\isaconst{invoke}\ \isavar{P}\ {\isacharparenleft}{\kern0pt}\isaconst{Suc}\ \isavar{n}{\isacharparenright}{\kern0pt}\ \isavar{st}\ \isavar{fd}\ \isavar{env}\ {\isacharequal}{\kern0pt}\isanewline
\isaconst{exec}\ \isavar{P}\ \isavar{n}\ {\isacharparenleft}{\kern0pt}\isaconst{fd{\isacharunderscore}{\kern0pt}ret}\ \isavar{fd}{\isacharparenright}{\kern0pt}\ \isavar{env}\ \isavar{st}\ {\isacharparenleft}{\kern0pt}\isaconst{fd{\isacharunderscore}{\kern0pt}body}\ \isavar{fd}{\isacharparenright}{\kern0pt}\ {\isasymbind}\isanewline
{\isacharparenleft}{\kern0pt}{\isasymlambda}{\isacharparenleft}{\kern0pt}\isabound{oc}{\isacharcomma}{\kern0pt}\ \isabound{env}{\isacharcomma}{\kern0pt}\ \isabound{st}{\isacharparenright}{\kern0pt}{\isachardot}{\kern0pt}\isanewline
\isaindent{{\isacharparenleft}{\kern0pt}\ \ \ }\isaconst{cleanup}\ \isavar{P}\ \isavar{n}\ \isabound{st}\ \isabound{env}\ {\isasymbind}\isanewline
\isaindent{{\isacharparenleft}{\kern0pt}\ \ \ }{\isacharparenleft}{\kern0pt}{\isasymlambda}\isabound{st}{\isachardot}{\kern0pt}\ \isaconst{Ok}\ {\isacharparenleft}{\kern0pt}case\ \isabound{oc}\ of\ \isaconst{Normal}\ {\isasymRightarrow}\ \isaconst{RUnit}\ {\isacharbar}{\kern0pt}\ \isaconst{Returned}\ \isabound{r}\ {\isasymRightarrow}\ \isabound{r}{\isacharcomma}{\kern0pt}\ \isabound{st}{\isacharparenright}{\kern0pt}{\isacharparenright}{\kern0pt}{\isacharparenright}{\kern0pt}%
\end{isabelle}%
\end{isamarkuptext}\isamarkuptrue%
%
}%endsnip
\snip{run_eqn}{1}{2}{%
%
\begin{isamarkuptext}%
%
\begin{isabelle}%
\isaconst{run}\ \isavar{fuel}\ \isavar{P}\ {\isacharequal}{\kern0pt}\isanewline
{\isacharparenleft}{\kern0pt}case\ \isaconst{find{\isacharunderscore}{\kern0pt}fun}\ \isavar{P}\ {\isacharprime}{\kern0pt}{\isacharprime}{\kern0pt}main{\isacharprime}{\kern0pt}{\isacharprime}{\kern0pt}\ of\ \isaconst{None}\ {\isasymRightarrow}\ \isaconst{Err}\ {\isacharparenleft}{\kern0pt}\isaconst{ErrUnknown}\ {\isacharprime}{\kern0pt}{\isacharprime}{\kern0pt}main{\isacharprime}{\kern0pt}{\isacharprime}{\kern0pt}{\isacharparenright}{\kern0pt}\isanewline
\isaindent{{\isacharparenleft}{\kern0pt}}{\isacharbar}{\kern0pt}\ \isaconst{Some}\ \isabound{fd}\ {\isasymRightarrow}\isanewline
\isaindent{{\isacharparenleft}{\kern0pt}{\isacharbar}{\kern0pt}\ \ \ }\isaconst{invoke}\ \isavar{P}\ \isavar{fuel}\ \isaconst{init{\isacharunderscore}{\kern0pt}state}\ \isabound{fd}\ {\isacharbrackleft}{\kern0pt}{\isacharbrackright}{\kern0pt}\ {\isasymbind}\ {\isacharparenleft}{\kern0pt}{\isasymlambda}{\isacharparenleft}{\kern0pt}\isabound{uu}{\isacharcomma}{\kern0pt}\ \isabound{st}{\isacharparenright}{\kern0pt}{\isachardot}{\kern0pt}\ \isaconst{Ok}\ {\isacharparenleft}{\kern0pt}\isaconst{out{\isacharunderscore}{\kern0pt}lines}\ \isabound{st}{\isacharparenright}{\kern0pt}{\isacharparenright}{\kern0pt}{\isacharparenright}{\kern0pt}%
\end{isabelle}%
\end{isamarkuptext}\isamarkuptrue%
%
}%endsnip


\lstdefinelanguage{Mojo}{
  morekeywords={def,struct,var,ref,mut,out,deinit,fn,if,else,while,return,print,for,raises,self,True,False,and,or,not,pass,from,import},
  sensitive=true, morecomment=[l]{\#}, morestring=[b]"}
\lstset{language=Mojo,basicstyle=\ttfamily\small,keywordstyle=\bfseries,
  commentstyle=\itshape\color{gray},columns=fullflexible,xleftmargin=1em,
  aboveskip=0.6em,belowskip=0.4em,showstringspaces=false}
\lstdefinestyle{out}{language={},basicstyle=\ttfamily\footnotesize}

\newenvironment{snippet}{\begin{quote}\small}{\end{quote}}
\newcommand{\micro}{micro-Mojo}
\newcommand{\IHOL}{Isabelle/HOL}
\newcommand{\code}[1]{\texttt{#1}}

\title{A Formal Semantics of Mojo 1.2 and Its Ownership Checker}
\author{Mark Utting\\ \small (with the assistance of Claude; see Acknowledgements)}
\date{\today\\[0.3em]\small Draft}

\begin{document}
\maketitle

\begin{abstract}
Mojo is a new systems programming language that combines Python-like syntax
with value ownership, argument conventions, first-class \emph{origins} and
as-soon-as-possible (ASAP) destruction of values. Its memory-safety rules are
described in prose and by the behaviour of a closed, rapidly evolving compiler;
there is no formal specification. We report on a formalisation in \IHOL{} of a
subset, \micro{}, that is small enough to be tractable but large enough to
exhibit the interesting ownership phenomena: a big-step semantics (a clocked
definitional interpreter with explicit errors), a flow-sensitive ownership and
type checker, an ASAP destructor-insertion pass, and a machine-checked proof
that checked programs never fail with an ownership or type error. A formal
model proves facts about \emph{itself}, so we also validate it against the real
compiler. We built a translator from Mojo source to the \IHOL{} syntax that
rejects programs outside the subset with an explanation, and used it to generate
an \IHOL{} theory that evaluates 76 Mojo programs together with the
behaviour (output, compile error or run-time error) that Mojo~1.2.0 showed. The
model agrees with the compiler on all of them, apart from two that expose a
documented gap; differential testing also exposed two further differences that
were closed, one by changing the model and one by narrowing the subset.
\end{abstract}

\section{Introduction}

Mojo~\cite{mojo-manual} aims to give Python programmers the performance of a
systems language. Unlike Python it has no garbage collector: every value has a
single owner, and the compiler decides when to destroy it. Unlike Rust, there
is no separate borrow-checker vocabulary visible in most code; instead,
functions declare how they receive each argument (\code{read}, \code{mut},
\code{var}, \code{ref}, \code{out}, \code{deinit}), references carry a static
\emph{origin}, and values are destroyed right after their last use rather than
at the end of a scope. These rules decide whether a program compiles, when
destructors run and therefore what a program prints, yet they are specified
only informally.

This paper describes a formalisation of that core in \IHOL{}~\cite{isabelle}.
Our contributions are:
\begin{enumerate}
\item a \emph{region model} of the memory-management proposal for Mojo
  origins~\cite{smith-gist}, with the design alternatives documented
  (Section~\ref{sec:region});
\item \micro{}: abstract syntax, an executable big-step semantics, an ownership
  and type checker, and an ASAP destructor-insertion pass (Section~\ref{sec:formal});
\item a machine-checked soundness theorem relating the checker to the semantics
  (Section~\ref{sec:sound});
\item a differential-testing pipeline from Mojo source to an \IHOL{} theory that is
  evaluated in batch mode, a suite of 76 programs, and an account of the
  differences it found (Sections~\ref{sec:dt} and~\ref{sec:findings}).
\end{enumerate}
The paper is also a case study in a way of working: the entire development,
including the proofs, the translator and the test programs, was produced by an
AI assistant under the direction of the author, with the proof assistant and
the Mojo compiler acting as the independent checkers (see the
Acknowledgements).

\section{Why formalise, and what makes Mojo hard}

\paragraph{Benefits.}
A formal semantics gives (i) \emph{precision}: writing a rule down forces a
decision about every corner the prose leaves open (for example: what happens to
a reference into a list when the list is appended to? when exactly is a
\code{var} parameter that the callee never uses destroyed?); (ii)
\emph{executability}: the semantics is a function, so example programs can be
run and their output compared with expectations before any theorem is proved;
(iii) \emph{guarantees}: soundness is a statement about all programs, not those
we thought to test; and (iv) a \emph{reference} against which a compiler or an
alternative language design can be compared. Because the model is written in a
proof assistant, claims such as ``the elaboration pass only inserts drops'' are
checked rather than argued.

\paragraph{Challenges specific to Mojo.}
\begin{itemize}
\item \emph{No specification.} The documentation describes intent; the compiler
  defines behaviour, and both are moving (we tested against Mojo~1.2.0; syntax
  such as \code{def} versus \code{fn} and the spelling of conventions has
  changed between releases). We treat the compiler as a black box oracle.
\item \emph{Destruction is observable and implicit.} Destructors (\code{\_\_deinit\_\_})
  can print, so \emph{when} a value is destroyed is part of a program's meaning.
  Mojo destroys values ASAP, after their last use, not at scope exit, so the
  semantics needs a last-use analysis, and what counts as a use includes uses of
  any reference whose origin contains the variable.
\item \emph{Many argument conventions.} Each convention changes what the caller
  may do afterwards (move, mutate, alias) and what the callee may do with the
  argument; implicit copies are legal for some types and not for others.
\item \emph{Origins are first-class.} A returned reference has a declared origin
  that is a set of parameters; checking a call must compute the origin of its
  result and extend the lifetimes of the variables in it.
\item \emph{Dynamic containers.} A reference into a list is invalidated by
  operations that change the list's shape but not by writes to its elements;
  the ownership tree of a list is not static.
\item \emph{Flow-sensitive checking.} Initialisation of variables and of
  individual fields (inside \code{\_\_init\_\_}) depends on control flow, so the
  checker is an abstract interpreter with joins at branches and a fixed point at
  loops.
\end{itemize}

\section{The region model}\label{sec:region}

The starting point is a design sketch by Nick Smith~\cite{smith-gist}, which
reinterprets an origin as a first-class, typed compile-time value denoting a
\emph{set of variables} a pointer may target. Our theory
\code{Mojo\_Region\_Model} gives that sketch a formal reading in four parts:
(1) a dereference-safety discipline that tracks initialisation of the variables
in a region; (2) dynamic containers whose item regions are invalidated precisely
when the container's \emph{shape} is mutated; (3) a replacement for Rust's
``aliasing XOR mutability'' by a disjointness rule on the region \emph{parameters}
of a function signature, so that pointers sharing a region parameter may alias
freely; and (4) mutability as an \emph{effect} declared by a function rather than
a property of an origin type. Every choice, and the alternatives not taken, is
discussed in comments in the theory. This model is deliberately abstract: it
speaks about regions and stores, not about programs. The later development
instantiates it, taking the ownership tree of the heap as the owner function.

\section{\micro{}: syntax, semantics and checker}\label{sec:formal}

\paragraph{The subset.}
\micro{} has \code{Int}, \code{Bool}, structs and \code{List[T]}; functions with
all six conventions; structs with fields, \code{@fieldwise\_init},
\code{Copyable}, \code{\_\_init\_\_(out self,\ldots)}, \code{\_\_deinit\_\_(deinit
self)} and methods; and the statements \code{var}, \code{ref} bindings,
assignment, \code{if}, \code{while}, \code{return}, \code{print} and an explicit
\code{drop}. Traits, generics, closures, \code{raises}, \code{for}, pointers and
overloading are left out (Section~\ref{sec:related} discusses extension).
Places (\code{x}, \code{p.f}, \code{p[i]}) are a separate syntactic category
from expressions, which lets the semantics decide from the callee's convention
whether an argument is passed as a location or as a value.

\subsection{Abstract Syntax}
The syntax is a set of \IHOL{} datatypes; the listings in this section are
generated from the theory sources, so they cannot drift from them. Types,
argument conventions and operators are enumerations:
\begin{snippet}
\begin{isabelle}
\usesnip{ast_ty}\\[0.3em]
\usesnip{ast_conv}\\[0.3em]
\usesnip{ast_binop}
\end{isabelle}
\end{snippet}
\code{TStr} is the type of string literals, which may only be printed, and
\code{TNone} the type of a call that returns nothing. \code{Imm} is Mojo's default
(read-only) convention. \code{Out} and \code{Deinit} are the conventions of the
\code{self} parameter of \code{\_\_init\_\_} and \code{\_\_deinit\_\_}. Places and
expressions are mutually recursive, because an index expression occurs inside a
place:
\begin{snippet}
\begin{isabelle}
\usesnip{ast_placeexpr}
\end{isabelle}
\end{snippet}
The comment on each constructor shows the Mojo syntax it represents. A \emph{place} names a location (\code{x}, \code{p.f}, \code{p[i]}). \code{EPlace p}
uses a place as an expression: read as a value it is an implicit copy (legal
only for trivial values), passed to a \code{mut}, \code{ref} or \code{imm}
parameter it is a reference. \code{ETake p} is Mojo's transfer \code{p\^{}}, and
\code{ECopy p} is \code{p.copy()}. A call \code{ECall f args} constructs a struct if
\code{f} names one and calls a function otherwise, as in Mojo. Statements are:
\begin{snippet}
\begin{isabelle}
\usesnip{ast_stmt}
\end{isabelle}
\end{snippet}
\code{SDrop} never appears in source programs; it is inserted by the ASAP
elaboration described below. Declarations follow:
\begin{snippet}
\begin{isabelle}
\usesnip{ast_ret}\\[0.3em]
\usesnip{ast_fundecl}\\[0.3em]
\usesnip{ast_structdecl}\\[0.3em]
\usesnip{ast_prog}
\end{isabelle}
\end{snippet}
A reference-returning function declares its result as \code{RRef [a, b] T},
Mojo's \code{-> ref[a, b] T}: the origin is the union of the listed parameters.

\subsection{Sugared Syntax}
Writing programs as raw constructor applications would be unreadable, so the
theory defines notation that mimics Mojo where Isabelle's syntax allows it and
uses upper-case keywords elsewhere (in the style of the IMP language). A variable
used as a place is marked with a section sign, a field access and an index look
like \code{p.f} and \code{p[i]}, and a place used as an expression is wrapped in
angle brackets. Calls are written with braces around the argument list, and
statements are sequenced with a double semicolon. Statements use the keywords
\code{VAR}, \code{REF}, \code{::=} (assignment), \code{+=}, \code{IF}, \code{WHILE},
\code{RETURN}, \code{PRINT}, \code{EVAL} (an expression statement) and \code{DROP};
\code{p += e} is sugar for assigning $p + e$ to $p$, as for Mojo's \code{Int}.
The sugar is purely notational: it is parsed to, and printed from, the datatypes
above, so every theorem about the syntax applies to sugared programs unchanged.

\paragraph{An example.}
All examples of the theory come as a Mojo program and its encoding. The first,
shown in Figure~\ref{fig:ex1}, exercises the conventions, a reference return, a
transfer and a list. It uses a resource whose destructor prints, so that the
order of destruction is visible in the output, and three small functions.
\begin{figure}[t]
\centering
\begin{minipage}[t]{0.5\textwidth}
\begin{lstlisting}[basicstyle=\ttfamily\footnotesize]
@fieldwise_init
struct Res:
    var id: Int
    def __deinit__(deinit self):
        print("drop", self.id)

def consume(var r: Res):
    print("got", r.id)

def bump(mut r: Res, n: Int):
    r.id += n

def pick(c: Bool, ref a: Res,
         ref b: Res) -> ref[a, b] Res:
    if c:
        return a
    else:
        return b
\end{lstlisting}
\end{minipage}%
\begin{minipage}[t]{0.47\textwidth}
\begin{lstlisting}[basicstyle=\ttfamily\footnotesize]
def main():
    var a = Res(1)
    var b = Res(2)
    bump(a, 10)
    ref r = pick(True, a, b)
    print(r.id)
    consume(a^)
    var l: List[Int] = [1, 2, 3]
    ref e = l[0]
    print(e)
    l.append(4)

# Output:
#   11
#   drop 2
#   got 11
#   drop 11
#   1
\end{lstlisting}
\end{minipage}
\caption{The first example: shared declarations (left) and \code{main} (right), with the
output that Mojo~1.2 produces.}
\label{fig:ex1}
\end{figure}
Mojo prints \code{11}, \code{drop 2}, \code{got 11}, \code{drop 11}, \code{1}. The
value \code{b} is destroyed straight after the last use of \code{r}, whose origin
is $\{a,b\}$, and before \code{consume} runs, while \code{a} is transferred to
\code{consume} and destroyed inside it. The final \code{append} invalidates
\code{e}, which is harmless because \code{e} is not used afterwards. The encoding of
\code{main} in the sugared syntax is:
\begin{snippet}
\begin{isabelle}
\usesnip{ex_ex1_main}
\end{isabelle}
\end{snippet}
Names are strings, \code{N 1} is the integer literal, and \code{TT} is \code{True}.
Note how the call \code{pick(True, a, b)} gets \code{\isasymlangle\isasymsection a\isasymrangle}
for its arguments: the callee's \code{ref} convention decides, at run time, that a
location and not a copy is passed.

\subsection{Memory Model}
Each variable owns a \emph{root} in a heap; a value is a tree (a struct holds its
fields, a list its items), and a location is a path from a root. Mojo's ownership
tree is thereby literal, a move transfers a subtree, and destroying a value
deinitialises a subtree. Uninitialised and moved-from locations hold a special
value, and touching one is the error \code{ErrUninit}. Every list carries a
\emph{generation} number that changes on shape-changing operations; a path
through a list records the generation it saw, so following a stale path is
\code{ErrInvalidated}. When a function returns, the roots of its locals are freed
and a reference into one is \code{ErrDangling}.

\subsection{Semantics}
The semantics is a total, fuel-bounded function in the style of functional
big-step semantics~\cite{owens-fbs}, with an error monad. Errors are explicit
results rather than stuck states, so the theorem ``checked programs never reach
these errors'' is directly expressible, and the semantics can be \emph{run}
with \code{value} and \code{eval}. Divergence appears as \code{OutOfFuel}. The
result type is \code{'a res = Ok 'a | Err err | OutOfFuel}, and the errors are
\code{ErrUninit}, \code{ErrDangling}, \code{ErrInvalidated}, \code{ErrBounds},
\code{ErrNotCopyable}, \code{ErrFieldInit} and the type and arity errors that the
checker rules out.

The interpreter is one group of 13 mutually recursive functions. Each takes a
fuel $n$ as its first argument after the fixed program $P$, and each recursive
call uses one unit of fuel, which makes termination trivial. Calls make
expressions depend on statements (a function body is a statement) and statements
on expressions, hence the single group:
\begin{center}\small
\begin{tabular}{@{}ll@{}}
\toprule
function & role\\
\midrule
\code{eval} & expression $\to$ \emph{owned} value\\
\code{evals} & list of expressions, left to right\\
\code{eval\_place} & place $\to$ path; evaluates and bounds-checks indices, and\\
 & records the list's generation in the path\\
\code{eval\_ref} & location-denoting expression (a place, or a call returning a\\
 & reference) $\to$ path\\
\code{pass\_args} & binds arguments to parameters by convention, giving the\\
 & callee's initial environment\\
\code{call\_fun}, \code{call\_meth} & functions, constructors, methods, and the built-in\\
 & list methods\\
\code{invoke} & runs a body in a fresh frame and destroys the frame's locals\\
\code{exec}, \code{exec\_block} & statements, and statements in a new block scope\\
\code{drop\_path}, \code{drop\_paths} & destroy a value: run \code{\_\_deinit\_\_}, then the\\
 & remaining fields or list items in order\\
\code{cleanup} & destroy and free the owned roots of an environment, newest first\\
\bottomrule
\end{tabular}
\end{center}
A value is a tree and a location a path, so \code{eval} returns an owned value
while \code{eval\_place} and \code{eval\_ref} return paths into the heap. An
environment maps each variable either to an owned root (\code{Own r}) or to a
borrowed path (\code{Bor q}); this is the only difference between a \code{var}
variable and a \code{mut} or \code{ref} parameter.

\code{pass\_args} is where the argument conventions are implemented. For
\code{var} (and \code{deinit}) the argument is evaluated to an owned value which
the callee owns in a fresh root, so \code{x\^{}} transfers and an implicit copy
is checked by the static rules to be of trivial type. For \code{mut} and
\code{ref} the argument must denote a location, which is borrowed. For
\code{imm} a location holding a trivial value is copied, since Mojo passes such
values by value, other locations are borrowed, and a non-location argument such
as a literal becomes a temporary owned by the callee. \code{call\_fun} first
looks the name up among the structs: for a struct with an \code{\_\_init\_\_} it
allocates a value whose fields are all uninitialised, passes it as \code{out self}
and afterwards checks that every field is set (\code{ErrFieldInit}); with
\code{@fieldwise\_init} it builds the value from the arguments; otherwise it
calls the function. The frame discipline lives in \code{invoke}, whose equation
is generated from the theory:
\begin{snippet}
\begin{isabelle}
\usesnip{invoke_eqn}
\end{isabelle}
\end{snippet}
Here \isa{m {\isasymbind} f} is bind in the error monad (written with \code{do}
notation in the theory): run the body, then \code{cleanup} the environment that
the body ended with, and return its result. \code{exec} threads an environment
and a heap through a statement and returns whether the statement ended normally
or with a \code{return}; \code{SVar} allocates a fresh root and extends the
environment, \code{SAssign} destroys the old value of the place before storing
the new one, and \code{SDrop p} calls \code{drop\_path}. \code{drop\_path} runs the
struct's \code{\_\_deinit\_\_} with \code{self} bound to the location and then
destroys the fields in order; this is the code that gives destruction its
observable effect, since a destructor can \code{print}. A program is run from
\code{main}:
\begin{snippet}
\begin{isabelle}
\usesnip{run_eqn}
\end{isabelle}
\end{snippet}
and the result is the list of printed lines (a line is a list of values). Output
printed before an error is not part of the result, because the error monad does
not carry the state.

\subsection{ASAP Elaboration}\label{sec:sound-intro}
To keep the dynamic semantics free of last-use analysis, a static pass
\code{elaborate} inserts explicit \code{DROP} statements: after each statement,
the owned variables it used that are not live afterwards are dropped, branches
drop variables that are live before the \code{if} but not within the branch, and
loop bodies are elaborated against the loop's live set. A use of a reference
counts as a use of every variable in its origin. Destroyed values run their
\code{\_\_deinit\_\_} and then the destructors of their remaining fields and list
items, in order. Programs without \code{DROP}s still have a C++-like meaning
(owned locals are destroyed in reverse order at scope end), which never fires
after elaboration.

\subsection{Checker}
\code{check\_prog} is a flow-sensitive abstract interpretation of each function
body, with an abstract state recording (i) which owned variables are definitely
initialised, (ii) which reference bindings may have been invalidated by a
shape change of a list they point into, and (iii) inside \code{\_\_init\_\_}, which
fields of \code{out self} are initialised. It also checks types, that implicit
copies are only of trivial values, that arguments are exclusive (comparing the
origins of borrowed arguments, the disjointness rule of the region model),
that immutable arguments are not mutated, and that a returned reference comes
from its declared origin. Some restrictions are deliberate simplifications
that make the checker \emph{stricter} than Mojo: variable names must be
distinct within a function, static paths stop at the first index (so
\code{l[i]} and \code{l[j]} always overlap), and analysis is at statement rather
than sub-expression granularity.

\section{Correctness results}\label{sec:sound}

The central invariant \code{inv} relates a static environment $\Gamma$ and
abstract state $\sigma$ to a run-time environment and heap: an owned variable
that is definitely initialised holds a well-typed value; a borrowed parameter
denotes a well-typed value (except the fields of \code{out self} not yet
initialised); and a reference binding denotes a path below its base variable
which is well-typed unless the reference is flagged as invalidated. Soundness
is proved by a single induction over the fuel, because expressions, statements
and calls are mutually recursive. The main theorem (here in \IHOL{} syntax,
slightly abridged) is:
\begin{lstlisting}[language={},basicstyle=\ttfamily\small]
theorem run_safe:
  assumes "wf_prog P"
  shows   "safe (run fuel P)"

corollary run_no_ownership_error:
  assumes "wf_prog P" and "run fuel P = Err e"
  shows   "e = ErrBounds" and "~ ownership_err e"
\end{lstlisting}
In words: running a program accepted by the checker, with any amount of fuel,
either terminates, runs out of fuel, or fails with an out-of-range index; it
never fails with a use of a moved or uninitialised value, a use of an invalid or
dangling reference, a copy of a non-copyable value, a return before \code{out
self} is fully initialised, or any type or arity error. The result is stronger
than ownership safety because the checker also checks types, and excluding
uninitialised reads needs well-typed values anyway. Out-of-bounds accesses are
not excluded, which matches Mojo, which reports them at run time. The same
statement is proved for every statement and expression (\code{stmt\_sound\_calls},
\code{expr\_sound\_calls}), and a separate lemma shows that elaboration only
inserts \code{DROP}s (\code{elab\_only\_inserts\_drops}), so safety carries over to
the elaborated program that is executed. The proofs occupy about 8\,500 lines of
\IHOL{}, over 340 lemmas, with the call case alone taking more than half.

\paragraph{What the theorem does not say.}
It says the \emph{model's} checker is sound for the \emph{model's} semantics. It
does not say the model resembles Mojo. A model that rejected every program would
satisfy it. This gap is the motivation for the next section.

\section{Differential testing}\label{sec:dt}

\paragraph{Idea.}
To validate fidelity, we generate Mojo programs, run them with the real compiler
to obtain the observed outcome, translate the same programs to \micro{}, and ask
\IHOL{} to evaluate them. Each test has one of three expected outcomes: the program
prints a given list of lines; it is rejected at compile time; or it fails at run
time. The \IHOL{} predicates are, respectively, that the elaborated program run by
the interpreter produces that output, that \code{check\_prog} reports an error,
and that the run ends with an error. Output is rendered as text (integers in
decimal, booleans as \code{True}/\code{False}, arguments space-separated) so
that it can be compared with the compiler's standard output.

\paragraph{Pipeline.}
\begin{center}
\small
\begin{tabular}{c}
\code{prog.mojo} $\xrightarrow{\text{translator}}$ \IHOL{} term of type \code{prog}\\[0.3em]
\code{prog.mojo} $\xrightarrow{\text{mojo build, run}}$ \code{prog.expected.json}\\[0.3em]
both $\xrightarrow{\text{generator}}$ \code{Micro\_Mojo\_Suite.thy}
$\xrightarrow{\text{isabelle build}}$ pass/fail
\end{tabular}
\end{center}
The translator is a hand-written tokenizer (with Python-style indentation), a
recursive-descent parser and a translator with type inference, in about 1\,500
lines of Python. Its most important property is what it \emph{refuses}. Any
construct outside the subset is rejected with a message that names the
feature and says why it is not supported, for example:
\begin{lstlisting}[style=out]
suite/rejected/a09_short_circuit.mojo:7:18: unsupported Mojo feature: right
operand of 'and' with calls or indexing
    Mojo short-circuits 'and'/'or' but the micro-Mojo semantics evaluates both
    operands; use nested 'if' statements when the right operand has side
    effects or can fail
\end{lstlisting}
A silent mistranslation would invalidate the whole experiment, so the
translator is conservative: it also rejects integer literals beyond $10^6$
(see below), and it is unit-tested on about forty rejection cases. We further
proved in \IHOL{} that the translator's output for the example programs of the
syntax theory is \emph{equal} to the hand-written encodings, so the two
independent encodings cross-check each other.

\paragraph{Recorded results.}
The driver records, for each test, the compiler's status, standard output and
diagnostics, and the compiler version, in a JSON file. The generated theory
therefore needs no Mojo installation to be re-evaluated: it is one \code{lemma}
per program, proved by \code{eval}, and the whole suite builds as an ordinary
\IHOL{} session in about two minutes.

\paragraph{The test programs.}
There are 11 programs from the original syntax theory and 65 new ones designed
to probe likely weak spots: timing of destruction (14), moves (including in
branches and loops), implicit and explicit copies, reference invalidation,
lists, \code{mut}/\code{var} parameters, methods and initialisers, and
arithmetic and control flow. Mojo accepted and ran 53 of the 76 programs,
rejected 20 at compile time and aborted 3 at run time (out-of-range index, pop
of an empty list, negative index). Compile-error tests are deliberately written
so that only one rule is violated, and we recorded Mojo's diagnostics to check
that the reason is the intended one (use of an uninitialised value,
invalidated interior reference, aliasing a mutably borrowed value, and so on).
We note one limitation: the \IHOL{} test asserts only \emph{that} the checker rejects
the program, not that it does so for the same reason.

The 11 original examples had expected outputs written by hand \emph{before} the
compiler was run; Mojo agreed with every one.

\paragraph{An example.}
Program \code{d13} has a struct whose destructor prints and which contains
three fields with destructors of their own:
\begin{lstlisting}
struct Outer:
    var a: Inner
    var b: Inner
    var c: Inner
    def __deinit__(deinit self):
        print("outer start")
        print("a is", self.a.id)
        print("outer mid")
        print("b is", self.b.id)
        print("outer end")
\end{lstlisting}
Mojo prints \code{drop inner 3} \emph{before} \code{outer start} (field \code{c} is
never used, so it is destroyed on entry), then destroys \code{a} immediately after
its last use and \code{b} likewise. A scope-based model, or one that runs the body
and then destroys the fields in order, gets a different trace.

\section{What differential testing found}\label{sec:findings}

Of the 65 new programs, 61 agreed on the first evaluation. The four that did not
revealed three kinds of gap, plus a tooling issue.

\begin{table}[t]
\centering
\begin{tabular}{@{}p{3.2cm}p{5.3cm}p{4.6cm}@{}}
\toprule
Difference & Mojo~1.2.0 & Resolution\\
\midrule
Unused \code{var} parameter (d09, d14) &
 destroyed when the callee is entered &
 \textbf{model fixed:} elaboration drops it on entry; soundness still holds\\[0.4em]
Fields in \code{\_\_deinit\_\_} (d05, d13) &
 each field destroyed ASAP; unused ones before the body &
 \textbf{documented divergence;} test asserts the mismatch\\[0.4em]
\code{and}/\code{or} (a09) &
 short-circuit &
 translator \textbf{rejects} impure right operands\\[0.4em]
Variable reuse in disjoint scopes (a08) &
 allowed &
 translator \textbf{rejects}; checker requires unique names\\
\bottomrule
\end{tabular}
\caption{Differences between the model and Mojo~1.2.0 found by the suite.}
\label{tab:diffs}
\end{table}

\paragraph{Unused \code{var} parameters.}
The model dropped an owned parameter only after its last use, so one that was
never used lived until the callee returned. Mojo destroys it on entry. The fix
was a small change to the elaboration pass: the body starts with drops of the
owned parameters that are not live on entry. Because the soundness proofs are
about the checker and interpreter, not the particular elaboration, the proofs
were unaffected, which we confirmed by rebuilding. This is the kind of change
the proof structure is meant to make cheap.

\paragraph{Field-wise destruction in \code{\_\_deinit\_\_}.}
In the model, \code{\_\_deinit\_\_} runs as a body over a borrowed \code{self}
and afterwards the remaining fields are destroyed in order. Mojo instead treats
the fields of \code{deinit self} as separately owned, and applies the ASAP rule
to each. Matching this would require tracking initialisation per field in the
checker's abstract state, a drop rule for fields of \code{deinit self}, an
elaboration over field places rather than variables, and changes to the
invariant used throughout the soundness proofs. We judged the cost not worth it
for a divergence that is visible only when a struct with printing field
destructors also has a printing \code{\_\_deinit\_\_}, and recorded it instead.
Each known divergence is a file entry with a reason, and the generated theory
asserts \emph{the negation} of the test for it. The build stays green today, and
becomes red as soon as someone fixes the semantics, at which point the entry is
deleted and the test becomes an ordinary one. This turns a known gap into a
regression test for the future fix.

\paragraph{Short-circuit \code{and}/\code{or}.}
The interpreter evaluates both operands. Making it short-circuit would change the
evaluation function and the induction in every soundness case that involves
boolean operators. Since the difference is observable only when the right-hand
operand has effects (a call that prints, or an indexing that may fail), we
narrowed the subset: the translator rejects such programs and explains why. The
rejected example is kept in a directory of its own, and a unit test checks that it
stays rejected, so that the restriction cannot silently disappear.

\paragraph{Variable names.}
Mojo allows the same \code{var} name to be declared in disjoint scopes. The
model's checker requires unique names in a function (this simplifies the
soundness invariant, since the environment can be treated as a finite map from
names to kinds). Again the translator rejects such programs.

\paragraph{A tooling observation.}
Evaluating integer arithmetic around $10^9$ with \code{eval} hung on the
(32-bit-pointer) Poly/ML platform we used, even for plain natural-number
arithmetic. The translator therefore limits literals to $10^6$ in absolute value.
We mention this because a differential test harness has its own trusted base,
and its failures can look like failures of the semantics.

\paragraph{Agreement on rejection.}
All 20 compile-error programs were rejected by the checker, and all three
run-time errors were reproduced as out-of-range accesses. In particular the
checker rejected partial moves, the use of references after a shape-changing
mutation, conflicting borrows of the same origin, and \code{\_\_init\_\_} that
does not initialise every field, each as Mojo did. Since the model's checker is
intentionally stricter than Mojo in places, the converse direction (programs
Mojo accepts but the model rejects) is the more informative one: the cases we
found there are the restrictions of Section~\ref{sec:formal}, which are surfaced
as translator rejections where they matter.

\section{Threats to validity}

\begin{itemize}
\item The suite was designed by the same agent that wrote the semantics, so it
  may share blind spots. The first-pass agreement on the hand-predicted
  examples is encouraging but not independent evidence. Random program
  generation \`a la Csmith~\cite{csmith} would complement it.
\item Tests assert the \emph{status} of a compile error, not its cause.
\item The suite targets one compiler version; Mojo changes quickly.
\item The theorem is relative to the model; each translator rejection and each
  recorded divergence is part of the model's stated scope and should be read
  alongside it.
\item Values are small (decimal integers below $10^6$); arithmetic overflow and
  floating point are not modelled.
\end{itemize}

\section{Related work}\label{sec:related}

\paragraph{Rust.} The closest bodies of work formalise Rust's ownership.
RustBelt~\cite{rustbelt} proves semantic soundness of a core calculus and its
unsafe libraries in Iris; Oxide~\cite{oxide} gives a syntactic type system for
borrowing based on ``provenances'' (regions of loans); Featherweight
Rust~\cite{pearce} models lifetimes and borrowing in a small calculus; Stacked
Borrows~\cite{stackedborrows} and Tree Borrows~\cite{treeborrows} define
aliasing models for unsafe code and are validated with the Miri interpreter.
Non-lexical lifetimes~\cite{rfc2094} are the analogue of our flow-sensitive
checker. Our work differs in the object language, which has argument
conventions, ASAP destruction and origins as sets of variables, and in method: a
single proof assistant development containing the semantics, the checker and a
compiler-validated test suite. Gao and Parreaux's InvalML~\cite{invalml}, raised in
the Mojo design discussion, is a closer model of invalidation.

\paragraph{Regions.} Tofte and Talpin's region calculus~\cite{tofte} and Cyclone's
regions~\cite{cyclone} are the origin of the idea that a pointer's type names
the memory it may refer to. Mojo's origins differ in that they are sets of
\emph{variables} and are first-class values.

\paragraph{Ownership in other languages.} Swift's \code{borrowing} and
\code{consuming} conventions~\cite{se0377} are similar to Mojo's argument
conventions; C++ destructors give scope-based rather than ASAP destruction.

\paragraph{Mechanised language semantics and testing.} CompCert~\cite{compcert},
the WebAssembly specification~\cite{watt}, JSCert~\cite{jscert} and the executable
C semantics~\cite{ellison} demonstrate the value of executable formal
semantics tested against implementations. Differential testing in general goes
back to McKeeman~\cite{mckeeman} and was made famous for compilers by
Csmith~\cite{csmith}. Our setting differs in that the oracle is a compiler for
a language whose semantics is not otherwise specified, and the object under test
is the formal semantics itself. The use of a clocked definitional interpreter
follows Owens et al.~\cite{owens-fbs}.

\section{Conclusion and future work}

We have shown that a core of Mojo's ownership discipline, including ASAP
destruction, argument conventions, origins and invalidation of references into
dynamic containers, can be given an executable semantics and a checker whose
soundness is proved in \IHOL{}, and that a pipeline from Mojo source to an
\IHOL{} theory lets the compiler itself be used as an oracle in batch mode. The
exercise produced concrete findings: a model bug fixed (unused \code{var}
parameters), a gap documented and kept as a regression test (field-wise
\code{\_\_deinit\_\_}), and two restrictions of the subset made explicit in the
translator. We found the combination valuable: the proof tells us that
the checker implies the absence of ownership errors, and the testing tells us
how far the model is from Mojo.

Future work: (i) field-wise destruction in \code{\_\_deinit\_\_}; (ii) short-circuit
operators and block scoping with shadowing; (iii) pointers with origins and
\code{for} loops over lists; (iv) traits, generics and \code{raises}; (v) random
generation of programs, including mutation of existing tests, and tracking the
causes of compile errors; (vi) re-running the suite on newer Mojo versions to
track the language's evolution.

\section*{Acknowledgements}

The author set the goals, took the design decisions that were put to them
(including the choice of a clocked functional big-step semantics and of an
explicit \code{drop} elaboration), directed each phase (region model,
subset proposal, syntax, semantics, checker, soundness, differential testing,
larger suite), asked for the compiler to be used once it was found to be
installed through \code{uv}, and decided what to commit and what to leave as a
documented divergence.

The development itself was carried out by Claude (Anthropic), working through
Claude Code, with the Isabelle PIDE MCP server for interactive proof
development; the final phases were done with Claude Sonnet~5.5. Claude wrote the
region theory and its commentary, the \micro{} syntax, semantics, checker and
elaboration, the soundness proofs, the Python translator and test driver, the
unit tests, the 76 Mojo test programs, the build configuration, the READMEs, the
analysis of each test disagreement, the small fix to the elaboration pass, and a
first draft of this paper. In the testing phase, Claude first wrote the expected output
for each example by hand from the semantics and only then ran the compiler, which
agreed. The validity of the formal results rests on the Isabelle kernel and
of the empirical results on the Mojo compiler, both independent of the
assistant that produced the artefacts; the statements of the definitions
and theorems, the choice of subset and this paper's claims should still be
read critically by a human, and the bibliography below has been written from
memory of the literature and needs checking against the sources.

\begin{thebibliography}{99}\small
\bibitem{mojo-manual} Modular. \emph{Mojo Manual: Ownership and value lifecycle}.
  \url{https://docs.modular.com/mojo/manual/}. Accessed 2026.
\bibitem{smith-gist} N.~Smith. \emph{An alternative model for ``lifetimes'' in Mojo}.
  GitHub gist and discussion.
  \url{https://gist.github.com/nmsmith/cdaa94aa74e8e0611221e65db8e41f7b}.
\bibitem{isabelle} T.~Nipkow, L.~C. Paulson, M.~Wenzel. \emph{Isabelle/HOL: A Proof
  Assistant for Higher-Order Logic}. LNCS 2283, Springer, 2002.
\bibitem{owens-fbs} S.~Owens, M.~O. Myreen, R.~Kumar, Y.~K. Tan. Functional big-step
  semantics. ESOP 2016.
\bibitem{rustbelt} R.~Jung, J.-H. Jourdan, R.~Krebbers, D.~Dreyer. RustBelt: Securing
  the foundations of the Rust programming language. POPL 2018.
\bibitem{oxide} A.~Weiss, D.~Patterson, N.~D. Matsakis, A.~Ahmed. Oxide: The essence of
  Rust. arXiv:1903.00982, 2019.
\bibitem{pearce} D.~J. Pearce. A lightweight formalism for reference lifetimes and
  borrowing in Rust. ACM TOPLAS 43(1), 2021.
\bibitem{stackedborrows} R.~Jung, H.-H. Dang, J.~Kang, D.~Dreyer. Stacked Borrows: An
  aliasing model for Rust. POPL 2020.
\bibitem{treeborrows} N.~Villani, J.~Hostert, D.~Dreyer, R.~Jung. Tree Borrows. PLDI 2025.
\bibitem{rfc2094} N.~D. Matsakis. Rust RFC 2094: Non-lexical lifetimes. 2017.
\bibitem{invalml} Gao and Parreaux. InvalML. OOPSLA 2025. (Reference as cited in the
  design discussion~\cite{smith-gist}; details to be verified.)
\bibitem{tofte} M.~Tofte, J.-P. Talpin. Region-based memory management. Information and
  Computation 132(2), 1997.
\bibitem{cyclone} D.~Grossman, G.~Morrisett, T.~Jim, M.~Hicks, Y.~Wang, J.~Cheney.
  Region-based memory management in Cyclone. PLDI 2002.
\bibitem{se0377} Swift Evolution SE-0377: \code{borrowing} and \code{consuming}
  parameter ownership modifiers.
\bibitem{compcert} X.~Leroy. Formal verification of a realistic compiler. CACM 52(7), 2009.
\bibitem{watt} C.~Watt. Mechanising and verifying the WebAssembly specification. CPP 2018.
\bibitem{jscert} M.~Bodin, A.~Charguéraud, D.~Filaretti, P.~Gardner, S.~Maffeis,
  D.~Naudziuniene, A.~Schmitt, G.~Smith. A trusted mechanised JavaScript
  specification. POPL 2014.
\bibitem{ellison} C.~Ellison, G.~Rosu. An executable formal semantics of C with
  applications. POPL 2012.
\bibitem{mckeeman} W.~M. McKeeman. Differential testing for software. Digital Technical
  Journal 10(1), 1998.
\bibitem{csmith} X.~Yang, Y.~Chen, E.~Eide, J.~Regehr. Finding and understanding bugs in
  C compilers. PLDI 2011.
\end{thebibliography}
\end{document}
